% BasicLayerReview.tex -- review of the foundational (Basic) layer of lean-DLMF as a
% possible starting point for incorporating DLMF into TauCeti and later Mathlib.
% Compile with LuaLaTeX: latexmk -lualatex BasicLayerReview.tex
\documentclass[11pt,a4paper]{article}
\usepackage{leanreport}
\hypersetup{pdftitle={lean-DLMF: review of the Basic layer, 2026-11-02}}

\newcommand{\dlmf}{\repofile{lean-DLMF}}
\newcommand{\resolved}[1]{\textit{Resolved #1.}}
\newcommand{\open}{\textit{Open.}}

\begin{document}

\title{The Basic Layer of \texttt{lean-DLMF}\\as a Starting Point for TauCeti and Mathlib}
\author{Bastiaan J. Braams}
\date{2026-11-02}
\maketitle

\begin{abstract}
  The plan is to point the Mathlib and TauCeti communities to the foundational files of
  \texttt{lean-DLMF} only: each \file{Basic.lean} and every file below a \file{Basic/}
  directory. Those files would then serve as the start for bringing all of DLMF into TauCeti
  first and Mathlib later. This review asks two questions. Do the Basic files contain the
  important definitions and their basic theorems? Do they import only other Basic files of
  \texttt{lean-DLMF}, besides Mathlib, \module{DLMFMathlib} and TauCeti? The first version of
  this review (2026-11-01) was of commit \texttt{a329483}. This version redoes it against the
  working tree of 2026-11-02 (commit \texttt{ca05ade} with uncommitted changes), after the
  reorganization that its recommendations started. Items settled since the first version are
  marked as resolved rather than removed. The supporting lists of the first version are in
  \dlmf{reports/BasicLayerInventory20261101.json}; they have not been regenerated. No library
  source was changed for this review.
\end{abstract}

\section{Summary}

The Basic layer has 235 files with 48{,}833 lines: 66 files \file{Basic.lean} (among them the
entry point \module{DLMF.Basic} and \module{DLMF.Core.Basic}) and 169 files below \file{Basic/}
directories. \module{DLMF.Basic} imports all of them. They hold 1{,}611 public definitions and
1{,}791 public theorems (2{,}004 counting private ones). They contain no \lean{sorry}, and all
1{,}770 DLMF-tagged declarations in them are proved. On 2026-11-01 the layer had 122 files,
15{,}265 lines, 1{,}233 public definitions and 180 public theorems.

\paragraph{Imports.} \resolved{2026-11-01} Every Basic file now imports only Mathlib, TauCeti,
\module{DLMFMathlib} and other Basic files, and the import closure of \module{DLMF.Basic}
contains no lean-codes module. \module{DLMF.Core.Framework} was renamed
\module{DLMF.Core.Basic}, \lean{symEllIntR} moved to the theorem layer, and the import survey
now fails on any violation of either rule (\S\ref{sec:imports}). Two small items remain: 13
Basic files import \module{DLMF.Core.Basic} without using it, and \module{DLMF.Core.Basic} has
not yet been split into \module{DLMFMathlib} and the areas.

\paragraph{Definitions.} The Basic layer now has most of the defining material: 1{,}128 of the
1{,}501 DLMF-tagged definitions in the compiled environment (878 of 1{,}503 on 2026-11-01).
\resolved{2026-11-01} Elementary, Constants, Zeta and NumberTheory have Basic layers, and almost
all of the central functions that the first version found in theorem files (Hankel, Whittaker,
\(Q_\nu\), the associated Legendre and Ferrers functions, Barnes' \(G\), the Hurwitz zeta
function, the polylogarithm, Lerch's \(\Phi\), the Legendre-form elliptic integrals, Anger and
Weber, the Askey scheme) have moved down unchanged. The overlaps of the first version are
settled. A few central definitions remain in theorem files, among them Weber's \(W(a,x)\), the
Riemann \(P\)-equation, the character Bernoulli numbers and \(q\)-factorial, and the new
namespace policy raises two questions for reviewers: namespaces that differ from the directory
(error functions in \lean{DLMF.Confluent}, \module{DLMF.Core.Basic} in \lean{DLMF}), and bridges
to TauCeti that sit in the theorem layer or are missing (\S\ref{sec:definitions}).

\paragraph{Basic theorems.} \resolved{in large part, 2026-11-01 and 2026-11-02} This was the
main gap. Of the 8{,}203 DLMF-tagged theorems, 642 are now in the Basic layer (16 on
2026-11-01). Of the 765 tagged theorems of the definitions sections \S\(n\).2, 353 are in Basic
files; for Chapters 4, 9 and 27 all of them are, for Chapters 32 and 34 all but the admitted
ones. The basic properties of Airy, Bessel, Gauss and Kummer functions, Legendre's
equation, the Weber equation, theta, Weierstrass and Heun functions are proved from foundational
imports. The rest of the gap is concentrated in Chapters 8, 11, 13, 14, 19, 22, 28, 29 and 33,
where the \S\(n\).2 results sit on theorem stacks of 20 to 120 modules, and in a dozen shallow
modules whose proved results could move almost unchanged (\S\ref{sec:theorems}).

\section{Imports}\label{sec:imports}

\subsection{Direct imports}

The 235 Basic files have these direct imports:

\begin{itemize}
  \item 498 of Mathlib modules (208 distinct), 15 of TauCeti modules (13 distinct) and 35 of
    \module{DLMFMathlib} modules (22 distinct);
  \item 437 of other Basic files.
\end{itemize}

None of them is a DLMF theorem module, a non-foundational \module{DLMF.Core} module or a
lean-codes module. \resolved{2026-11-01} The 14 imports of \module{DLMF.Core.Framework} became
imports of the foundational \module{DLMF.Core.Basic}, and the two Carlson imports of
\file{DLMF/EllipticIntegrals/Basic/Symmetric.lean} are gone. The four legacy files
\file{Definitions.lean} (RiemannTheta, OrthogonalPolynomials, JacobiElliptic, Theta) lie below
\file{Basic/} directories, so they are Basic files by location.

\subsection{The closure of \texorpdfstring{\texttt{DLMF.Basic}}{DLMF.Basic}}

\module{DLMF.Basic} imports 74 modules (61 on 2026-11-01), and its transitive closure reaches
all 234 other Basic files. Inside the project, the closure contains no module that is not a
Basic file. It also contains 26 \module{DLMFMathlib} modules (13 on 2026-11-01), among them
\module{DLMFMathlib.NumberTheory.PrimeNumberTheorem}, which imports TauCeti's prime number
theorem, and \module{DLMFMathlib.NumberTheory.ModularForms.KleinJSurjective}, the home of
Klein's \(J\). Its modules import 235 Mathlib modules and 15 TauCeti modules directly.
\resolved{2026-11-01} The closure contained 183 lean-codes modules, all through the Carlson
imports; it now contains none.

The import survey, run as
\begin{center}
  \texttt{python3 scripts/import\_survey.py --baseline scripts/import\_baseline.json},
\end{center}
reports no foundational or closure violations, no missing project modules and no file-import
cycles. Its one new item is the two-way area pair Elementary \(\leftrightarrow\)
IncompleteGamma, between the theorem modules \module{DLMF.Elementary.LogExpIntegrals},
\module{DLMF.IncompleteGamma.FractionTools} and \module{DLMF.IncompleteGamma.SpecialValues}. Of the 21 two-way
area pairs, five run through a Basic file on one side (Confluent \(\leftrightarrow\) Coulomb,
Confluent \(\leftrightarrow\) ParabolicCylinder, Confluent \(\leftrightarrow\)
OrthogonalPolynomials, EllipticIntegrals \(\leftrightarrow\) JacobiElliptic and Polylog
\(\leftrightarrow\) Zeta); since areas have no hierarchy, they are not violations.

\subsection{\texorpdfstring{\texttt{DLMF.Core.Basic}}{DLMF.Core.Basic}}

\resolved{2026-11-01} \file{DLMF/Core/Framework.lean} was renamed \file{DLMF/Core/Basic.lean}
with no declaration changes, and its 36 importers were updated. It now belongs to the
foundational layer.

The file (232 lines) fixes shared conventions. It defines improper integrals
(\lean{HasIntegralAtTop}, \lean{improperAtTop} and siblings), the principal value
\lean{HasPrincipalValueIoi}, straight-segment integrals and primitives (\lean{segmentIntegral},
\lean{primitive}), the sector filters \lean{sector} and \lean{logSector}, finite continued
fractions (\lean{cfContinuants}, \lean{cfTail}, \lean{HasContinuedFraction}),
\lean{oddProduct}, the spherical Bessel functions \lean{sphericalJ} and \lean{sphericalI},
\lean{orderedZero} and \lean{asymptoticCoeff}.

\open{} Of the 17 Basic files that import it (besides \module{DLMF.Basic}), only four use one
of its definitions: \file{ErrorFunctions/Basic/Functions.lean} (\lean{primitive}),
\file{IncompleteGamma/Basic.lean} and \file{GeneralizedHypergeometric/Basic/MeijerG.lean}
(\lean{segmentIntegral}), and \file{Spheroidal/Basic/EquationAndSpectralData.lean}
(\lean{sphericalJ}). The other 13 import it only for the Mathlib modules that it imports in
turn. The trimming recommended in the first version was not done; it is harmless now that
\module{DLMF.Core.Basic} is foundational, but it hides the real dependencies.

\open{} The second step of the first version still applies: give each definition its natural
home, the general analysis (improper integrals, principal values, segment primitives, sector
filters, continued fractions, \lean{asymptoticCoeff}) in \module{DLMFMathlib} in the namespaces
of the Mathlib APIs it extends, \lean{sphericalJ}, \lean{sphericalI} and \lean{oddProduct} in
\file{Bessel/Basic.lean} (DLMF \S10.47), and \lean{orderedZero} in AlgebraicAnalytic or
AsymptoticApproximations. The namespace part is resolved later on 2026-11-02:
\lean{sphericalJ} and \lean{sphericalI} moved to \lean{DLMF.Bessel} in
\file{Bessel/Basic.lean}, and the remaining declarations of \module{DLMF.Core.Basic} are in
\lean{DLMF.Core}. The remaining Core modules
(\module{MathlibCoverage}, \module{Upstream}, \module{UpstreamElementary},
\module{ZeroCounting}) are not imported by any Basic file.

\subsection{The Carlson imports}

\resolved{2026-11-01} \lean{symEllIntR} (DLMF 19.16.8), the only definition in a Basic file
built on Carlson's \(R\)-function, moved unchanged into
\file{DLMF/EllipticIntegrals/Symmetric.lean}. \file{EllipticIntegrals/Basic/Symmetric.lean} now
imports Mathlib only and keeps the integral definitions \lean{symEllIntRF}, \lean{symEllIntRJ},
\lean{symEllIntRG}, \lean{symEllIntRD}, \lean{symEllIntRC} and \lean{symEllIntS}. The survey
fails if a lean-codes module re-enters the closure of \module{DLMF.Basic}.

Outside the Basic layer, lean-codes still sits low in many theorem stacks. Most theorem modules
that hold \S\(n\).2 results reach 188 to 226 lean-codes modules, and they do so through three
Hypergeometric files:
\begin{itemize}
  \item \module{DLMF.Hypergeometric.EulerIntegral}, through its import of Carlson's
    \module{Carlson.TwoVariable.GaussHypergeometric};
  \item \module{DLMF.Hypergeometric.Transformations};
  \item \module{DLMF.Hypergeometric.FurtherTransformations.Cubic}.
\end{itemize}
This matters for the moves of
\S\ref{sec:theorems}: a property module that does not need those three files can usually be
detached from lean-codes entirely.

\section{Definitions}\label{sec:definitions}

\subsection{Areas without a Basic file}

\resolved{2026-11-01} Elementary (DLMF Chapter 4), Zeta (Chapter 25), NumberTheory (Chapter 27)
and Constants now have Basic layers, with nine, five, four and one files. They hold the
reciprocal trigonometric and hyperbolic functions, the inverse functions with their DLMF
branches, the Gudermannian, the general logarithm and power values, and, since 2026-11-02, the
Lambert \(W\) branches with their explicit, Mathlib-only construction
(\file{Elementary/Basic/LambertWConstruction.lean} and \file{Elementary/Basic/LambertW.lean});
\(\xi\), \(\chi\), the Stieltjes constants, Riemann--Siegel \(\theta\) and \(Z\), the complex
Hurwitz zeta function, the polylogarithm, Lerch's \(\Phi\) and the Fermi--Dirac and
Bose--Einstein integrals; Jordan's totient, Ramanujan and Gauss sums, \lean{ramanujanTau},
Lambert series and the sum-of-squares counts; and Catalan's constant.

One area is left without a Basic file: WaveEquation, whose six definitions (two- and
three-dimensional Laplacians, cylindrical and spherical waves, heat modes) illustrate
applications and have no DLMF chapter of their own. That is acceptable. The top-level file
\file{DLMF/Bailey.lean} (\lean{IsBaileyPair}, DLMF \S17.12) also holds a definition outside the
Basic layer, next to \lean{IsBaileyChain} in \file{QHypergeometric/BaileyPairs.lean}; the two
belong together in \file{QHypergeometric/Basic}.

AlgebraicAnalytic (Chapter 1) and AsymptoticApproximations (Chapter 2) have no area-level
\file{Basic.lean}, but they have eleven and seven \file{Basic.lean} files in their
subdirectories, which hold 183 and 113 definitions. Their definitions are well covered; their
theorems are not (\S\ref{sec:theorems}).

\subsection{Central definitions in theorem files}

Of the 1{,}501 DLMF-tagged definitions, 373 are outside the Basic layer (625 on 2026-11-01).
Most of them are auxiliary data and should stay where they are: coefficients of asymptotic and
uniform expansions (\lean{pcfOlverA}, \lean{ellAsymA}, \lean{gammaUniformS}), error-bound
multipliers, the reduction data of \S19.29 (\lean{SymmetricReduction.quarticProduct}, \lean{quadraticPairU}),
argument-reduction data for computation (\lean{expReductionIndex}), continued-fraction data,
zero-counting functions, and the predicates for conjectures (\lean{GoldbachConjecture},
\lean{LehmerConjecture}). \lean{symEllIntR} stays in the theorem layer by design.

\resolved{2026-11-01} Of the central functions listed in the first version, all but four now
have their home in the Basic layer: Barnes' \(G\) and Glaisher--Kinkelin, the real digamma and
\lean{scaledGamma}; the inverse error functions, \lean{iteratedErfc} and the Voigt functions;
the real and complex incomplete beta functions; the Scorer functions and the Airy modulus and
phase; the Hankel functions; Anger, Weber and Lommel functions; Whittaker's \(M\) and \(W\);
\lean{legendreQ}, the associated Legendre, Ferrers, Olver and conical functions and
\lean{sphericalHarmonic}; \lean{gaussOperator}; \lean{meijerGLog}; the \(q\)-exponential and
\(q\)-trigonometric functions and the \(q\)-Appell functions; \lean{chebyshevV},
\lean{chebyshevW}, the Askey scheme and the \(q\)-families, the Bessel and Pollaczek
polynomials, \lean{IsOrthogonalPolynomials}, \lean{christoffelDarbouxKernel} and
\lean{hankelDeterminant}; the Legendre-form and Bulirsch integrals; the inverse Jacobian
functions, Jacobi's \(\varepsilon\) and \(Z\); \lean{modularLambda}, \lean{dedekindSum},
\lean{etaMultiplier}, \lean{IsModularForm} and \lean{IsModularFunction}; the higher-order
Bernoulli and Euler polynomials, Genocchi and tangent numbers; and the combinatorial numbers of
Chapter 26. The four that remain are \lean{genExpIntL} (\S8.19), \lean{riemannOperator}
(\S15.11), \lean{gen3F2UnitThomae} (\S16.4) and \lean{IsBaileyPair}.

\open{} A second pass over the 373 definitions finds these central or near-central functions
still in theorem files:
\begin{itemize}
  \item Chapter 5: \lean{qFactorial} (5.18.2) and \lean{qBeta} (5.18.11) in
    \file{Gamma/QGamma.lean}, while \lean{qGamma} is in \file{QSeries/Basic/Pochhammer.lean};
    \lean{binetPhi}.
  \item Chapter 12: Weber's \(W(a,x)\) (\lean{pcfW}, \S12.14, in
    \file{ParabolicCylinder/Weber/Solutions.lean}), the second standard pair of solutions of
    Weber's equation; its file sits on 96 theorem modules.
  \item Chapter 15: \lean{riemannOperator} and \lean{riemannSolutions} (\S15.11), and the
    logarithmic second solution \lean{logarithmicSecond} of \S15.10 for integer \(c\).
  \item Chapter 16: \lean{genNormalizedODE} and \lean{genD} (\S16.8), the differential
    equation of \({}_{q+1}F_q\).
  \item Chapter 11: \lean{angerA} (\S11.10), the companion of \lean{angerJ} and
    \lean{weberE}.
  \item Chapter 18: the orthogonal polynomials on the unit circle
    (\lean{IsUnitCircleOP}, \lean{verblunskyCoefficient}, \S18.33) and the bivariate disk and
    triangle polynomials (\S18.37).
  \item Chapter 24: \lean{bernoulliChar} and \lean{bernoulliCharPoly} (24.16.10--11),
    \lean{degenerateBernoulli} and \lean{bernoulliSecondKind} (\S24.16), the siblings of the
    higher-order polynomials that moved down.
  \item Chapter 21: \lean{IsSymplectic} and \lean{symplecticJ} (\S21.5).
\end{itemize}
Most have short bodies; \lean{qFactorial} and \lean{qBeta} depend on a theorem module with no
theorem module below it, and \lean{bernoulliChar} only on Basic material.

\subsection{Overlaps}

The overlaps of the first version are settled:
\begin{itemize}
  \item Polylogarithm. \resolved{2026-11-01} \lean{DLMF.Zeta.polylog} is the canonical
    polylogarithm and \lean{DLMF.Zeta.polylogSeries} the canonical series 25.12.10; the
    natural-order \lean{Polylog.polylogSeries} is removed. The real integral forms
    \lean{Polylog.Real.dilog} and \lean{Polylog.Real.trilog} remain, bridged by
    \lean{Zeta.polylog\_two\_ofReal\_eq\_dilog}.
  \item Spherical harmonics and Jacobi polynomials. \resolved{2026-11-01} Not duplicates; the
    docstrings now cross-reference \file{SphericalHarmonics/Basic.lean} and
    \lean{Legendre.sphericalHarmonic}, and \lean{MatrixArgument.jacobiP} and
    \lean{OrthogonalPolynomials.jacobiPolyP}.
  \item Klein's \(J\). \resolved{2026-11-01} \lean{ModularForm.kleinJ} in
    \module{DLMFMathlib} is the one definition; \lean{Weierstrass.kleinJ} is removed and
    23.15.7 is the theorem \lean{Weierstrass.kleinJ\_eq\_jacobiTheta}.
  \item Complete elliptic integrals. \resolved{2026-11-01} \lean{compEllIntKk} and
    \lean{compEllIntKpk} moved to \file{EllipticIntegrals/Basic/Complete.lean}, beside their
    siblings.
  \item \lean{arctanCutPlane} and \lean{arcsinhDomain}: \resolved{2026-11-01} by the
    \lean{rfl} theorem \lean{arcsinhDomain\_eq\_arctanCutPlane}. \lean{cbinom}:
    \resolved{2026-11-01} removed in favour of Mathlib's \lean{Ring.choose}.
  \item Hurwitz zeta. \resolved{2026-11-01} \lean{hurwitzZetaComplex} moved from
    \lean{DLMF.Confluent} to \lean{DLMF.Zeta}, with the Basic bridge
    \lean{hurwitzZetaComplex\_eq\_hurwitzZeta} to Mathlib's function.
\end{itemize}

\open{} Checked against the upstream survey \dlmf{reports/UpstreamNamespaces.tex}, there
are two new overlaps with TauCeti:
\begin{itemize}
  \item The Hermite function \(\varphi_n\) of 1.18.42, defined complex-valued as
    \lean{hermiteFunction} in \file{AlgebraicAnalytic/Spectral/Basic.lean}, equals the
    coercion of TauCeti's real \lean{TauCeti.hermiteFunction}: both are \(\mathit{He}_n(\sqrt2 x)\,e^{-x^2/2}
    /\sqrt{n!\sqrt\pi}\). There is no bridge.
  \item The character Bernoulli numbers \(B_{n,\chi}\) of 24.16.10, defined as
    \lean{bernoulliChar} in the theorem layer, are TauCeti's
    \lean{DirichletCharacter.generalizedBernoulli}, with the same closed form
    \(f^{n-1}\sum_a\chi(a)B_n(a/f)\). There is no bridge.
\end{itemize}

\subsection{Namespaces and upstream bridges}

The namespace policy adopted on 2026-11-02 puts every declaration under \file{DLMF/} in
\lean{DLMF.<Area>}, real-variable forms in the \lean{Real} sub-namespace of the area, and asks
for a foundational bridge wherever an upstream function differs from DLMF's.

The Basic layer complies with the first part: apart from the dot-notation lemma
\lean{DirichletCharacter.isDirichletCharacterFn}, which the policy allows, no public
declaration of a Basic file lies outside \lean{DLMF}. \resolved{2026-11-02} The Gamma area moved
from Mathlib's \lean{Real} and \lean{Complex} namespaces into \lean{DLMF.Gamma} and
\lean{DLMF.Gamma.Real}, so that \lean{DLMF.Gamma.Real.digamma} now stands beside Mathlib's
\lean{Complex.digamma}, with the bridge \lean{digamma\_eq\_re\_digamma}.

\open{} A reviewer will notice that the namespace does not always name the directory. 22 Basic
files declare, wholly or partly, into another area: the error functions, exponential integrals,
incomplete gamma functions, parabolic cylinder functions, the Coulomb functions and the helper
\lean{hermitePolyHEval} of \file{OrthogonalPolynomials/Basic/Definitions.lean} live in
\lean{DLMF.Confluent} (the arrangement of Temme's book); the generalized hypergeometric and
Appell files in \lean{DLMF.Hypergeometric}; the complex incomplete beta function of
\file{IncompleteGamma/Basic/IncompleteBeta.lean} in \lean{DLMF.Statistical}; and
\file{Elliptic/Basic/ThetaPeriodicity.lean} in \lean{DLMF.JacobiElliptic}. Each is a
legitimate choice, but a contribution should either align namespace and directory or state the
rule. \emph{Decided later on 2026-11-02:} the Confluent and Hypergeometric namespaces stay, and
\file{AGENTS.md} states that an area namespace may span several directories; the declarations
of \module{DLMF.Core.Basic} moved to \lean{DLMF.Core}.

\begin{RaggedRight}
\resolved{2026-11-02} New foundational bridges: \lean{DLMF.Bessel.besselI\_ofReal} and
\lean{re\_besselI\_ofReal} to TauCeti's \lean{Real.besselI}, beside the existing
\lean{DLMF.Confluent.erf\_ofReal} (TauCeti's \lean{TauCeti.Real.erf}),
\lean{DLMF.Statistical.normIncBetaI\_ofReal} (\lean{TauCeti.regularizedIncompleteBeta}) and the
thin wrapper \lean{DLMF.Confluent.Real.invErf} of \lean{Real.erfInv}.

\open{} Bridges that the policy places in the foundational layer but that are still in theorem
files:
\begin{itemize}
  \item \lean{DLMF.MatrixArgument.multivariateGamma\_ofReal} (to
    \lean{TauCeti.multivariateGamma}), in \file{MatrixArgument/GammaBeta/GammaIntegrals.lean};
    its proof only unfolds both definitions. \emph{Resolved later on 2026-11-02:} it moved to
    \file{MatrixArgument/Basic.lean} unchanged, with \lean{multivariateGamma\_one}.
  \item \lean{DLMF.Confluent.lowerIncGamma\_ofReal} (to \lean{TauCeti.lowerIncompleteGamma}) in
    \file{Confluent/IncompleteGamma.lean}, and \lean{normIncGammaP\_ofReal},
    \lean{normIncGammaQ\_ofReal} (to \lean{TauCeti.regularizedGamma}) in
    \file{Statistical/IncompleteGamma.lean}. Their proofs use the integral 8.2.1, so they wait
    for the \S8.2 results to become foundational.
  \item \lean{Real.logIntLi\_eq\_logIntLi\_two\_add\_logIntegral} (to
    \lean{TauCeti.Real.logIntegral}) in \file{ExpIntegrals/AnalyticProperties.lean}; it needs
    the derivative of \(\operatorname{li}\) from foundational imports.
\end{itemize}
\end{RaggedRight}

\subsection{Definitions a Mathlib reviewer will question}

\begin{itemize}
  \item Parameter restrictions. \open{} \lean{laguerrePolyL}, \lean{jacobiPolyP} and
    \lean{ultrasphPoly} are still defined through \({}_1F_1\) and \({}_2F_1\) with a
    Pochhammer prefactor, ``for \(\alpha+1\) (resp.\ \(\lambda+\tfrac12\)) not a nonpositive
    integer'', so at the excluded parameters the value is not the polynomial. Mathlib would
    expect a definition valid for every parameter, preferably as a \lean{Polynomial} with an
    evaluation lemma, as Mathlib does for \lean{Polynomial.hermite} and
    \lean{Polynomial.Chebyshev.T}. The hypergeometric form then becomes a theorem.
  \item Limits in the order. \resolved{in part, 2026-11-01} \lean{besselY} and \lean{besselK}
    still define integer orders by \lean{limUnder} in the order parameter, but the Basic layer
    now proves what makes that usable: joint continuity in order and argument
    (\lean{continuousAt\_besselY}, \lean{continuousAt\_besselY\_order} and the \(K\)
    analogues), the recurrences and the Wronskians for every order. Analyticity in \(\nu\) is
    not yet stated.
  \item Real-only definitions. \resolved{2026-11-01 and 2026-11-02} The policy now fixes the
    choice: a real-only form lives in the \lean{Real} sub-namespace of its area, beside a
    complex form or an upstream function, with a bridge. \lean{DLMF.Gamma.Real.digamma}
    is bridged to \lean{Complex.digamma}, \lean{Polylog.Real.dilog} to
    \lean{Zeta.polylog}, and \lean{Statistical.Real.incBeta} has the complex
    \lean{Statistical.incBeta} beside it.
  \item Definitions by choice. \open{} Mathieu, Lam\'e and spheroidal characteristic values
    and functions are defined by choice from spectral data. Their existence theorems are in
    the theorem layer: \lean{mathieuSpecA\_exists}, \lean{mathieuSpecB\_exists},
    \lean{mathieuA\_spec} and \lean{mathieuB\_spec} in \module{DLMF.Mathieu.Basic.SpectralExistence},
    which is fully proved and has no theorem module below it, so it can move to
    \file{Mathieu/Basic} unchanged; \lean{exists\_spheroidalEigenvalue} and
    \lean{existsUnique\_isSpheroidalPsFactor} in \module{DLMF.Spheroidal.Existence} (four
    theorem modules below it); and \lean{Real.lameEigenvalue\_spec} (\S29.3(i)) in
    \module{DLMF.Lame.Eigenvalues}, which sits on 52 theorem modules and 189 lean-codes
    modules.
\end{itemize}

\section{Basic theorems}\label{sec:theorems}

\subsection{What is there}

The 235 Basic files contain 1{,}791 public theorems, and 642 tagged theorem entries. 110 files
contain definitions only, which matches the layout that keeps definition files short and puts
properties in separate files below \file{Basic/}. The largest theorem files are
\file{Bessel/Basic/SecondKind.lean} (62 public theorems), \file{Airy/Basic/Properties.lean}
(59), \file{Elementary/Basic/LambertWConstruction.lean} (52),
\file{Hypergeometric/Basic/ParameterAnalyticity.lean} (44),
\file{CoalescingSaddles/Basic/ContoursAndSymmetries.lean} and
\file{Mathieu/Basic/Properties.lean} (38 each).

\resolved{2026-11-02} The existence and analyticity of the principal continuation of
\({}_{q+1}F_q\) for all complex parameters are foundational
(\file{GeneralizedHypergeometric/Basic/PrincipalContinuationExistence.lean}), together with the
generalized-series modules it needs, which moved to \file{Hypergeometric/Basic/}. The
definition \lean{genRegularizedHG} therefore no longer depends on a theorem-layer existence
proof.

\open{} Areas whose Basic layer has definitions but almost no theorems: IncompleteGamma (8
definitions, no theorems), MatrixArgument (34, none), Selberg (19, none), Lam\'e (61, 2),
Coulomb (51, 9), Spheroidal (52, 13), and the methods chapters AlgebraicAnalytic (183, 3) and
AsymptoticApproximations (113, 1). For the last two most definitions are notions (contours,
distributions, asymptotic scales) whose theory is the chapter itself.

Table~\ref{tab:n2} shows, chapter by chapter, how many of the tagged theorems of \S\(n\).2 are
in Basic files. On 2026-11-01 the answer was none for every chapter except 5, 20 and 23.

\begin{table}[ht]
  \centering
  \caption{Tagged theorems of \S\(n\).2 in the Basic layer, of all tagged theorems of
    \S\(n\).2 (tag dump of 2026-11-02).}\label{tab:n2}
  \begin{tabularx}{\textwidth}{@{}l>{\RaggedRight\arraybackslash}X@{}}
    \toprule
    Coverage & Chapters (in Basic / total) \\
    \midrule
    all & 4 (21/21), 9 (18/18), 27 (16/16) \\
    all proved ones & 32 (40/47), 34 (4/5) \\
    half or more & 5 (3/5), 10 (4/6), 15 (5/10), 16 (4/6), 17 (28/54), 18 (12/22),
      20 (26/28), 21 (16/21), 23 (20/28), 24 (11/16), 30 (6/9), 31 (17/24) \\
    less than half & 6 (10/22), 12 (12/27), 22 (12/27), 25 (5/14), 28 (33/78),
      33 (6/23), 36 (11/30) \\
    a few & 7 (2/9), 13 (3/42), 14 (3/16), 19 (4/36), 29 (1/14) \\
    none & 1 (0/42), 2 (0/11), 8 (0/20), 11 (0/15), 35 (0/3) \\
    \midrule
    total & 353/765 \\
    \bottomrule
  \end{tabularx}
\end{table}

Where the remaining \S\(n\).2 results are:
\begin{RaggedRight}
\begin{itemize}
  \item Chapter 8 (20): \file{IncompleteGamma/ContinuedProperties.lean} and
    \file{Confluent/IncompleteGamma.lean}, \file{Statistical/IncompleteGamma.lean};
  \item Chapter 11 (15): \file{Struve/SeriesAndEquations.lean};
  \item Chapter 13 (39): \file{Confluent/BasicProperties/}, mainly
    \file{IntegerParametersAndLimits.lean}, \file{Wronskians.lean} and
    \file{Continuity.lean}, and \file{Confluent/ConnectionWronskians.lean};
  \item Chapter 14 (13): \file{Legendre/Associated/} (Wronskians, boundary values and
    connections, conjugation);
  \item Chapter 19 (32): \file{EllipticIntegrals/LegendreIdentities.lean}, with 25 proved;
  \item Chapter 22 (15): \file{JacobiElliptic/AnalyticProperties.lean} and the special
    values;
  \item Chapter 28 (45): \file{Mathieu/Properties/} (spectra, eigenfunctions, Floquet
    theory) and \file{Mathieu/SpectralExistence.lean};
  \item Chapter 29 (13): \file{Lame/DifferentialEquations.lean};
  \item Chapter 33 (17): \file{Coulomb/RhoEtaProperties.lean} and the Coulomb files of
    Confluent;
  \item Chapters 1 and 2 (53): elementary algebra (\file{AlgebraicAnalytic/Algebra/}) and
    \file{AsymptoticApproximations/Basics/Transcendental.lean}.
\end{itemize}
\end{RaggedRight}
Project-wide, 1{,}669 of the tagged theorems are admitted, but only 35 of the 405 \S\(n\).2
theorems outside the Basic layer are (the count includes three in \module{DLMFMathlib}); almost all of the gap is proved material in the wrong layer.

\subsection{How deep the key properties sit}

How hard it is to move a property into the Basic layer depends on what its proof imports. For
the modules that hold the remaining \S\(n\).2 results, the number of DLMF theorem modules in the
import closure, and of lean-codes modules, is:
\begin{RaggedRight}
\begin{itemize}
  \item shallow (proved results that can move almost unchanged):
    \begin{itemize}
      \item \module{DLMF.Mathieu.Basic.SpectralExistence}: 0 theorem modules, 0 lean-codes modules;
      \item \module{DLMF.Weierstrass.EllipticFunctionsAndInvariants}: 0 and 0 (6 proved);
      \item \module{DLMF.RiemannTheta.Scaled} and \module{Characteristics}: 0--1 and 0;
      \item \module{DLMF.QSeries.Pochhammer}, \module{QBinomial}, \module{Saalschutz},
        \module{QExponential}: 0--5 and 0;
      \item \module{DLMF.OrthogonalPolynomials.General} and \module{GeneralIdentities}: 0--1
        and 0;
      \item \module{DLMF.Zeta.Riemann}: 1 and 0 (6 proved, 1 admitted);
      \item \module{DLMF.Spheroidal.Existence}: 4 and 0;
      \item \module{DLMF.Bailey}: 5 and 0;
    \end{itemize}
  \item moderate:
    \begin{itemize}
      \item \module{DLMF.Hypergeometric.AnalyticProperties}: 17 and 188;
      \item \module{DLMF.Mathieu.Properties.Spectra}, \module{Eigenfunctions},
        \module{Floquet}: 19--22 and 27;
      \item the modules \module{Continuity}, \module{Wronskians} and
        \module{IntegerParametersAndLimits} of \module{DLMF.Confluent.BasicProperties}:
        22--24 and 188--222;
      \item \module{DLMF.Legendre.Associated.Wronskians}: 35 and 188;
    \end{itemize}
  \item deep:
    \begin{itemize}
      \item \module{DLMF.JacobiElliptic.AnalyticProperties}: 60 and 189;
      \item \module{DLMF.EllipticIntegrals.LegendreIdentities}: 66 and 192;
      \item \module{DLMF.Struve.SeriesAndEquations}: 79 and 222;
      \item \module{DLMF.Lame.DifferentialEquations}: 90 and 189;
      \item \module{DLMF.Coulomb.RhoEtaProperties}: 115 and 226;
      \item \module{DLMF.ExpIntegrals.AnalyticProperties}: 119 and 226;
      \item \module{DLMF.IncompleteGamma.ContinuedProperties}: 121 and 226;
      \item \module{DLMF.Spheroidal.Equation}: 132 and 226;
      \item \module{DLMF.CoalescingSaddles.CanonicalIntegrals.OscillatoryRepresentations}: 154
        and 226.
    \end{itemize}
\end{itemize}
\end{RaggedRight}
\resolved{2026-11-01} The worst case of the first version is gone: the Airy equation,
Wronskians and connection formulas are now proved in \file{Airy/Basic/Properties.lean}, whose
closure contains no theorem module, while the old \module{DLMF.Airy.Equation} still sits on 137.
\resolved{2026-11-01} Likewise the shallow cases named in the first version moved (Painlev\'e
\S32.2 except the admitted results, theta \S20.2, most of the Riemann zeta \S25.2 and of the
general orthogonal-polynomial identities of \S18.2). The deep cases are again mostly
accidental: a property module imports a large sibling (an asymptotics or error-bound framework,
or the Carlson-based Hypergeometric files) for a few lemmas.

\subsection{Recommendation}

\begin{enumerate}
  \item \resolved{2026-11-01} The project rule is refined as recommended. AGENTS.md now lists
    the content of the foundational layer for each central function (definition with
    docstring, unfolding and special values, symmetries, continuity or analyticity, the
    defining equation and Wronskian, the upstream bridge), requires that it be proved from
    foundational imports with no admitted dependency, and keeps everything else in the theorem
    layer.
  \item \resolved{2026-11-01} Properties go in files below \file{Basic/}, for example
    \file{Airy/Basic/Properties.lean}, and the survey covers them.
  \item \open{} Move the shallow cases above, starting with
    \module{DLMF.Mathieu.Basic.SpectralExistence} (which also answers the reviewer's question about
    definitions by choice) and \module{DLMF.Weierstrass.EllipticFunctionsAndInvariants}.
  \item \open{} For the deep cases, prove the key properties again from the Basic definitions
    instead of moving the existing proofs, as was done for Airy, Bessel, Gauss and Kummer. The
    most valuable remaining families are the Kummer Wronskians and integer-parameter limits
    (\S13.2), the incomplete gamma functions (\S8.2, which also unblocks two TauCeti bridges),
    the Legendre Wronskians (\S14.2), the Legendre-form identities (\S19.2) and the Struve
    series (\S11.2).
  \item \resolved{2026-11-01} Module docstrings in Mathlib style: every Basic file with five or
    more public theorems has a ``Main results'' section.
\end{enumerate}

\section{Recommended order of work}

Done since 2026-11-01: the import changes (rename of Framework, \lean{symEllIntR}, survey
checks); Basic files for Elementary, Constants, Zeta and NumberTheory and the moves of the
central definitions; the overlaps of the first version; the key-properties layer for most core
families; Lambert \(W\) and the \({}_{q+1}F_q\) continuation in the foundational layer; the
Gamma namespace move and the Bessel \(I\) bridge. What remains, in order:

\begin{enumerate}
  \item Shallow moves, one build: \module{DLMF.Mathieu.Basic.SpectralExistence},
    \module{DLMF.Weierstrass.EllipticFunctionsAndInvariants}, the proved \S\(n\).2 results of
    the shallow QSeries, RiemannTheta, OrthogonalPolynomials and Zeta modules
    (\lean{multivariateGamma\_ofReal} is done).
  \item The remaining central definitions of \S\ref{sec:definitions}, moved unchanged where
    the body allows (\lean{qFactorial}, \lean{qBeta}, \lean{bernoulliChar},
    \lean{IsBaileyPair}, \lean{riemannOperator}, \lean{genNormalizedODE}), and new bridges
    for \lean{hermiteFunction} and \lean{bernoulliChar}.
  \item The deep \S\(n\).2 families, area by area, starting with Chapters 13 and 8; then the
    incomplete gamma bridges.
  \item The parameter-restricted polynomial definitions of Chapter 18.
  \item Namespaces: align namespace and directory (or record the rule), and split
    \module{DLMF.Core.Basic} into \module{DLMFMathlib} and the areas. Both are renames and
    need the usual review. Trim the 13 unused imports of \module{DLMF.Core.Basic} at the same
    time.
\end{enumerate}

\section*{Method}

Foundational files were found by name, as in \dlmf{scripts/import_survey.py}: a module named
\file{Basic}, a module below a \file{Basic/} directory, or a legacy
\file{Defs}/\file{Definitions} module; the count of 235 includes \module{DLMF.Basic} and
\module{DLMF.Core.Basic}. Imports were read from the sources, and closures were computed over
the project and the sibling \file{../lean-codes} checkout; the import survey was run with its
baseline. Counts of tagged declarations come from a dump of the compiled environment of
2026-11-02 (\dlmf{scripts/DumpDLMF.lean}), with 10{,}744 entries, of which 9{,}704 belong to
\texttt{lean-DLMF}; a theorem counts for \S\(n\).2 when its tag lies in \S\(n\).2. Counts of
public declarations come from a source scan (top-level \texttt{def}, \texttt{abbrev},
\texttt{structure}, \texttt{class}, \texttt{inductive}, \texttt{theorem} and \texttt{lemma} not
marked \texttt{private}), and namespaces from the \texttt{namespace} commands of each file.
Whether a definition counts as ``central'' is the reviewer's judgment, made against the DLMF
chapter structure.

\end{document}
